% Generated draft: requires mathematical and source review.
\documentclass[11pt]{article}
\usepackage[margin=25mm]{geometry}
\usepackage{fontspec}
\setmainfont{Latin Modern Roman}
\usepackage{amsmath,amssymb,mathtools,unicode-math}
\setmathfont{Latin Modern Math}
\usepackage{enumitem,needspace,xurl,hyperref}
\setlength{\parindent}{0pt}
\setlength{\parskip}{5pt}
\setlength{\emergencystretch}{3em}
\newcommand{\N}{\mathbb{N}}
\newcommand{\Z}{\mathbb{Z}}
\newcommand{\Q}{\mathbb{Q}}
\newcommand{\R}{\mathbb{R}}
\newcommand{\C}{\mathbb{C}}
\newcommand{\F}{\mathbb{F}}
\newcommand{\A}{\mathbb{A}}
\newcommand{\PP}{\mathbb{P}}
\newcommand{\E}{\mathbb{E}}
\newcommand{\eps}{\varepsilon}
\newcommand{\dd}{\,\mathrm{d}}
\begin{document}
\section*{UnsolvedMath ID: 700975}
\textbf{Source problem code:} catalogue entry
\subsection{Definitions and mathematical statement}
Let \(p\) be a prime, let \(G\) be a finite group, and let \(\varphi\in\operatorname{Aut}(G)\) have order exactly \(p\). Define
\[
C_G(\varphi)=\{g\in G:\varphi(g)=g\},\qquad
Z(G)=\{z\in G:zg=gz\text{ for every }g\in G\}.
\]
Throughout, assume that \(C_G(\varphi)\leq Z(G)\). Define the derived series by \(G^{(0)}=G\) and
\[
G^{(n+1)}=[G^{(n)},G^{(n)}]\qquad(n\geq0),
\]
where \([H,H]\) is the subgroup generated by the commutators \(x^{-1}y^{-1}xy\) with \(x,y\in H\). The group \(G\) is soluble if \(G^{(n)}=1\) for some \(n\geq0\); its derived length \(d(G)\) is the least such \(n\).

The source asks the following questions.
\begin{enumerate}
\item[(a)] If \(p=3\), must \(G\) be soluble?
\item[(b)] Does there exist a function \(f\), from the set of primes to the nonnegative integers, such that for every prime \(p\), every finite soluble group \(G\), and every automorphism \(\varphi\) of \(G\) of order \(p\) satisfying \(C_G(\varphi)\leq Z(G)\), one has
\[
d(G)\leq f(p)?
\]
\item[(c)] As asked by V. K. Kharchenko, does there exist a function \(h\), from the set of primes to the nonnegative integers, such that for every prime \(p\), every finite \(p\)-group \(G\), and every automorphism \(\varphi\) of \(G\) of order \(p\) satisfying \(C_G(\varphi)\leq Z(G)\), one has
\[
d(G)\leq h(p)?
\]
Here a finite \(p\)-group has order \(p^m\) for some nonnegative integer \(m\), and the prime is the same as the order of \(\varphi\).
\end{enumerate}

The supplied source reports that V. D. Mazurov and T. L. Nedorezov proved solubility under the standing hypotheses when \(p=2\), and that nonsoluble examples exist for every prime \(p>3\). For part (c), it attributes to V. V. Bludov a simple example showing that nilpotency class cannot be bounded in terms of \(p\). Here the nilpotency class is the least \(c\geq0\) such that \(\gamma_{c+1}(G)=1\), where \(\gamma_1(G)=G\) and \(\gamma_{i+1}(G)=[\gamma_i(G),G]\). These background assertions are preserved as source reports; the record does not supply the examples or a separate citation for Bludov's construction.
\subsection{Short English statement}
For a finite group with an automorphism of prime order \(p\) whose fixed elements are central, ask whether the group must be soluble when \(p=3\). Also ask whether its derived length, when it is soluble, is bounded solely by \(p\), and whether such a bound holds in the special case of a finite \(p\)-group. The supplied registry labels the problem open; this status and the reported background results have not been independently verified.
\subsection{Sources}
\begin{itemize}
\item Local UnsolvedMath registry, supplied complete record with permanent ID \(700975\), exact source code \texttt{catalogue entry}, title \emph{Kourovka Notebook Problem 13.57}; upstream revision \texttt{d371ee917ac9647fbb06cd6fc529deec6f3dda82}. The record was created and updated on \texttt{2026-10-16T14:00:00Z} and carries the supplied status \texttt{open}. These are registry metadata, not independent mathematical verification.
\item \emph{The Kourovka Notebook (Unsolved Problems in Group Theory)}, original problem 13.57, 13th issue (1995); also identified in the supplied background with the 21st issue edition (2026), \texttt{arXiv:1401.0300v48}. Proposed by E. I. Khukhro. The supplied collection description names E. I. Khukhro and V. D. Mazurov as editors.
\item V. D. Mazurov and T. L. Nedorezov, \emph{Algebra and Logic} \textbf{35}, no. 6 (1996), 392--397. The supplied statement cites this reference for solubility when \(p=2\); no article title is supplied.
\item Part (c) is attributed to V. K. Kharchenko. The nilpotency-class example mentioned there is attributed to V. V. Bludov, without a separate bibliographic reference.
\end{itemize}
\end{document}
